%% ============================================================================
%%  math_commands.tex  —  custom math macros
%%
%%  Layout (each block sorted alphabetically by command name):
%%    (1) \let ..........  base symbols (a single letter or command) and
%%                         primitive aliases
%%    (2) \newcommand ...  aesthetic transformations built on the symbols above,
%%                         plus short upright-text tags
%%    (3) mathtools .....  \DeclareMathOperator, \DeclarePairedDelimiter and
%%                         \NewDocumentCommand definitions
%% ============================================================================


%% ---------------------------------------------------------------------------
%%  (1) \let : base symbols and primitive aliases
%% ---------------------------------------------------------------------------
\let\adjacencymatrixsymb A
\let\bigOnotationsymb O
\let\continuousfieldsymb \psi
\let\couplingmatrixsymb J
\let\critsymb c
\let\de \partial
\let\degreematsymb D
\let\densitymatrixsymb \rho
\let\diracdeltasymb \delta
\let\distancematrixsymbol D
\let\edgesetsymb E
\let\entropysymb S
\let\graphsignsymb \sigma
\let\graphsymb G
\let\hamiltoniansymb H
\let\Im \relax            % clear amsmath's \Im so it can be redeclared in (3)
\let\laplaciansymb L
\let\nodedegreesymb k
\let\partitionfunctionsymb Z
\let\probabilitysymb P
\let\propagatorsymbol K
\let\quenchavg \overline
\let\signedgraphsymb \Gamma
\let\sinkoperatorsymb S
\let\spinsymb \sigma
\let\transp \intercal
\let\vec \bm
\let\vertexsymb V
\let\weightedadjacencymatrixsymb W


%% ---------------------------------------------------------------------------
%%  (2) \newcommand : aesthetic transformations and upright-text tags
%% ---------------------------------------------------------------------------
\newcommand{\adj}{\adjacencymatrixsymb}
\newcommand{\Adjacency}{\wadj}                  % adjacency matrix (alias for \wadj)
\newcommand{\Band}{\mathcal{B}}                 % frequency band
\newcommand{\cfield}{\continuousfieldsymb}
\newcommand{\crit}{{\rm \critsymb}}
\newcommand{\Dcoph}{\tilde{\distm}}             % cophenetic communication distance matrix
\newcommand{\ddelta}{\diracdeltasymb}
\newcommand{\degm}{\degreematsymb}
\newcommand{\degreeshrt}{deg}
\newcommand{\diagonalmatrixshrt}{diag}
\newcommand{\dil}{{\rm \dilutedshrt}}
\newcommand{\dilutedshrt}{dil}
\newcommand{\distm}{\mathcal{\distancematrixsymbol}}
\newcommand{\dm}{\hat{\densitymatrixsymb}}
\newcommand{\edges}{\edgesetsymb}
\newcommand{\eff}{{\rm \effectiveshrt}}
\newcommand{\effectiveshrt}{eff}
\newcommand{\ent}{\entropysymb}
\newcommand{\graph}{\mathcal{\graphsymb}}
\newcommand{\ham}{\mathscr{\hamiltoniansymb}}
\newcommand{\highgamma}{\gamma_{\rm h}}
\newcommand{\imcoh}{\(|\mathrm{ImCoh}|\)}
\newcommand{\jij}{\couplingmatrixsymb}
\newcommand{\lapl}{\laplaciansymb}
\newcommand{\lowgamma}{\gamma_{\rm l}}
\newcommand{\MSC}{\Coherence}                   % alias for the \Coherence operator in (3)
\newcommand{\ndegree}{\nodedegreesymb}
\newcommand{\ord}{\bigOnotationsymb}
\newcommand{\partf}{\partitionfunctionsymb}
\newcommand{\Phase}{\Phi}
\newcommand{\pmassemp}{{p_{\rm mass}}}
\newcommand{\prob}{\probabilitysymb}
\newcommand{\propag}{\propagatorsymbol}
\newcommand{\rew}{{\rm \rewiringshrt}}
\newcommand{\rewiringshrt}{rew}
\newcommand{\rhocoph}{{\rho}^{\rm coph}}        % per-pair multiscale shift correlation
\newcommand{\rhodrift}{{\rho}_{\rm drift}}      % drift-floor null
\newcommand{\rhoraw}{{\rho}^{\rm raw}}
\newcommand{\rhoxprobe}{{\rho}_{\rm xp}}        % cross-probe restriction
\newcommand{\sdegm}{\bar{\degm}}
\newcommand{\setmea}{\abs}                      % alias for the \abs delimiter in (3)
\newcommand{\signshrt}{sign}
\newcommand{\sink}{{\rm \sinkshrt}}
\newcommand{\sinkop}{\sinkoperatorsymb}
\newcommand{\sinkshrt}{snk}
\newcommand{\Slapl}{\bar{\lapl}}
\newcommand{\SoftMask}{\mathbf{G}}              % soft mask
\newcommand{\spin}{\spinsymb}
\newcommand{\Surrogates}{\wadj^{\mathrm{null}}}
\newcommand{\TimeSeries}{\mathbf{X}}            % multichannel timeseries matrix
\newcommand{\TimeSeriesSingle}{x}               % single-channel signal
\newcommand{\traceshrt}{Tr}
\newcommand{\txtacr}[1]{\text{\acrshort{#1}}}
\newcommand{\vertexs}{\vertexsymb}
\newcommand{\wadj}{\weightedadjacencymatrixsymb}


%% ---------------------------------------------------------------------------
%%  (3) mathtools : operators, paired delimiters, document commands
%% ---------------------------------------------------------------------------

% --- math operators (alphabetical) ---
\DeclareMathOperator{\Coh}{Coh}
\DeclareMathOperator{\Coherence}{MSC}
\DeclareMathOperator{\corr}{corr}
\DeclareMathOperator{\degree}{\diagonalmatrixshrt}
\DeclareMathOperator{\diag}{\diagonalmatrixshrt}
\DeclareMathOperator{\Im}{Im}
\DeclareMathOperator{\ImCoh}{ImCoh}
\DeclareMathOperator{\Lapl}{\nabla^2}
\DeclareMathOperator{\mthdef}{\;\overset{\mathrm{def}}{=}\;}
\DeclareMathOperator{\rank}{rank}
\DeclareMathOperator{\sign}{\signshrt}
\DeclareMathOperator{\Tr}{Tr}

% --- paired delimiters (alphabetical) ---
\DeclarePairedDelimiter{\abs}{|}{|}
\DeclarePairedDelimiter{\avg}{\langle}{\rangle}
\DeclarePairedDelimiter{\bra}{\langle}{\rvert}
\DeclarePairedDelimiterX{\braket}[2]{\langle}{\rangle}{#1\,\delimsize\vert\,\mathopen{}#2}
\DeclarePairedDelimiter{\ket}{|}{\rangle}
\DeclarePairedDelimiter{\norm}{\lVert}{\rVert}
\DeclarePairedDelimiter{\qua}{[}{]}

% --- document commands: differentials and derivatives (alphabetical) ---
% Define differential operator
\NewDocumentCommand{\dd}{s o m}{
  \IfBooleanTF{#1} % Check if * is used
    {\mathrm{d}#3} % No space for inline form
    {\IfNoValueTF{#2} % Check if power is given
      {\mathrm{d}#3} % Normal differential
      {\mathrm{d}^{#2}#3} % Differential with power
    }
}

% Define total derivative
\NewDocumentCommand{\dv}{s o m m}{
  \IfBooleanTF{#1} % Check if * is used (inline form)
    {{\mathrm{d}#3}/{\mathrm{d}#4}} % Inline derivative
    {\IfNoValueTF{#2} % If no order is given
      {\frac{\mathrm{d}#3}{\mathrm{d}#4}} % First derivative
      {\frac{\mathrm{d}^{#2}#3}{\mathrm{d}#4^{#2}}} % Higher-order derivative
    }
}

% Define partial derivative
\NewDocumentCommand{\pdv}{s o m m}{
  \IfBooleanTF{#1} % Check if * is used (inline form)
    {{\partial #3}/{\partial #4}} % Inline partial derivative
    {\IfNoValueTF{#2} % If no order is given
      {\frac{\partial #3}{\partial #4}} % First partial derivative
      {\frac{\partial^{#2} #3}{\partial #4^{#2}}} % Higher-order partial derivative
    }
}